% 03-physical.tex - Physical System Interfaces

\chapter{Physical System Interfaces}
\label{sec:physical}

\section{Mechanical Interface}

\begin{newInVersion}{2.4.1}
The DR consists of a 1U rack-mount server (e.g., Gigabyte R152-Z30) with
internal hot-swap drive bays for storage.  No specific server model is required;
any x86\_64 system with sufficient CPU cores, RAM, and drive bays is suitable.
The server is mounted in a standard EIA 19'' rack.
\end{newInVersion}

\section{Network Interfaces}

The DR server connects to the station network with two Ethernet interfaces:

\begin{longtable}{l l p{7.5cm}}
\toprule
\textbf{Interface} & \textbf{Speed} & \textbf{Purpose} \\
\midrule
\endfirsthead
\endhead
\bottomrule
\endfoot

Data interface    & 10~GbE (or faster)  &
    Receives data frames from NDP.  Connected to the NDP data network.
    The UDP port is configured by the \texttt{DataInPort} parameter
    in \texttt{defaults\_v2.cfg}. \\

Control interface & 1~GbE              &
    Sends and receives MCS command/response messages.  Connected to the
    station control network.  The UDP ports are configured by the
    \texttt{MessageInPort} and \texttt{MessageOutPort} parameters
    in \texttt{defaults\_v2.cfg}. \\

\end{longtable}

\section{Storage}

\begin{newInVersion}{2.4.1}
The DR supports two storage configurations.  Single-drive storage is preferred
for new deployments; RAID-0 arrays remain supported for legacy DRSU
compatibility.
\end{newInVersion}

\subsection{Preferred: Single Drives}

Individual hard drives are installed in internal hot-swap bays and formatted
with the EXT4 file system.  Each drive is mounted under the
\texttt{/LWA\_STORAGE/} hierarchy and managed independently by DROS2.
This configuration simplifies drive replacement and eliminates the risk of
array-wide data loss from a single drive failure.

\subsection{Supported: RAID-0 Arrays (Legacy DRSU)}

For compatibility with existing Data Recorder Storage Units (DRSUs), DROS2
also supports RAID-0 arrays constructed with \texttt{mdadm}.  A DRSU
typically consists of multiple drives combined into a single \texttt{/dev/md}
device, formatted with EXT4, and mounted under \texttt{/LWA\_STORAGE/}.

RAID-0 arrays provide higher sequential write bandwidth but offer no
redundancy: failure of any constituent drive results in loss of all data on
the array.

\subsection{Storage Hierarchy}

Regardless of configuration, all storage is mounted under
\texttt{/LWA\_STORAGE/} with the following structure:

\begin{lstlisting}
/LWA_STORAGE/
  Internal/
    0/DROS/Rec/       # Raw recordings
    0/DROS/Spec/      # Spectrometer output
  External/
    0/DROS/           # External storage (copy/dump targets)
\end{lstlisting}

In a multi-DR deployment (Section~\ref{sec:multi-dr}), each instance uses a
separate storage root, e.g., \texttt{/LWA\_STORAGE/DR1/},
\texttt{/LWA\_STORAGE/DR2/}.

Storage setup (partitioning, formatting, and mounting) is managed by the
\texttt{StorageControl.sh} script included with DROS2.  DROS2 itself writes
to mounted file systems and does not perform low-level storage management.
